\chapter*{约定、记号与单位}
\addcontentsline{toc}{chapter}{约定、记号与单位}

如未注明，本书所进行的操作均是在实数域下进行的。

\begin{enumerate}[label=(\arabic*)]
\item \textbf{公式。}本文中的矩阵均用大写黑正体表示，如
\[
\mathbf A,
\]
而在表达其矩阵元时，会使用括号加下标或者小写字母加下标的形式，如
\[
(\mathbf A)_{ij},\qquad a_{ij}.
\]
向量通常用小写黑斜体表示，如
\[
\bm x,
\]
向量分量则表示为
\[
x_i.
\]
单位矩阵记作 $\mathbf I$：
\[
(\mathbf I)_{ij}=\delta_{ij},
\]
而当需要明确指定其阶数 $n$ 时，记作
\[
\mathbf I_n.
\]

\item \textbf{伪代码。}
\begin{enumerate}[label=(\roman*)]
\item 循环结构：
\begin{lstlisting}[style=bellacode,numbers=none]
do i = 1,n
    ...
end

while(cond)
    ...
end

do
    ...
end while(cond)
\end{lstlisting}
\item 判断结构：
\begin{lstlisting}[style=bellacode,numbers=none]
if(cond)
    ...
elseif(cond)
    ...
else
    ...
end
\end{lstlisting}
\item 二元运算：
\begin{itemize}
\item \texttt{a \&\& b}（与）
\item \texttt{a || b}（或）
\item \texttt{a.b}（指标缩并：$\sum_j a_{ij}b_{jk}$）
\end{itemize}
\end{enumerate}
\end{enumerate}
\clearpage
